\section{Conclusion and Future Work}
\label{sec:conclusion}

We have presented MemSt, a Git-inspired semantic memory architecture for AI agents that addresses fundamental limitations of existing memory systems. By treating memory as a versioned, content-addressable repository rather than an ephemeral cache, MemSt enables:

\begin{enumerate}
    \item \textbf{Hierarchical Retrieval}: Multi-level organization (session summaries $\rightarrow$ key turns $\rightarrow$ full context) with dynamic token budget allocation, achieving \textbf{41.7\% F1} on LOCOMO---an 84\% improvement over RAG baselines.
    
    \item \textbf{Specialized Graph Architectures}: TemporalKG achieves 54.0\% F1 on temporal questions (+120\% vs RAG), while MultiHopKG achieves 52.8\% F1 on multi-hop questions (+133\% vs RAG).
    
    \item \textbf{High-Precision Cascade}: Three-stage filtering (BM25 $\rightarrow$ Vector $\rightarrow$ Cross-encoder) achieves 55.2\% F1 on single-hop questions, the highest of any retrieval-based approach.
    
    \item \textbf{Exceptional Efficiency}: 3.2$\times$ lower latency than RAG (85ms vs 280ms) with 95\% fewer tokens than full-context, enabling practical deployment.
    
    \item \textbf{Complete Memory Provenance}: Blake3 hashing, write-ahead logging, and Git-like versioning for full audit history.
    
    \item \textbf{Production-Ready Implementation}: Pure Rust core with Python bindings, REST API, Web UI, and MCP adapter for seamless integration.
\end{enumerate}

\subsection{Key Insights}

Our comprehensive evaluation on LOCOMO (1,986 questions across 10 long conversations) reveals several key insights for memory system design:

\textbf{Hierarchy Beats Flat Retrieval.} Hierarchical Summary's strong performance across categories (49.5\% single-hop, 47.2\% multi-hop, 42.8\% temporal) validates the importance of structured memory organization that mirrors human cognitive processes.

\textbf{Specialization Matters.} The 16\% gap between best single-hop (Cascade) and best multi-hop (MultiHopKG) strategies suggests that question-aware routing could achieve higher overall performance. Our QuestionAdaptive strategy demonstrates this potential.

\textbf{Latency-Accuracy Tradeoffs Are Tunable.} Cascade offers highest accuracy at 180ms, while Hierarchical offers best efficiency at 85ms. Users can select appropriate tradeoffs for their application requirements.

\textbf{Category-Specific Search Weights Significantly Impact Accuracy.} Our ablation studies show that optimal BM25/vector weights vary by question type: 70:30 for single-hop (keyword-heavy) versus 30:70 for open-domain (semantic-heavy).

\subsection{Future Directions}

Several research directions remain:

\textbf{Question-Aware Routing.} Implementing a learned router to select optimal retrieval strategies based on question classification, potentially achieving 50\%+ F1 by combining specialized strengths.

\textbf{Learned Memory Policies.} While MemSt uses heuristic policies for tier transitions, reinforcement learning could learn optimal promotion/eviction strategies for specific domains.

\textbf{Cross-Encoder Integration.} Replacing our heuristic cross-encoder approximation with a trained neural re-ranker for the Cascade strategy.

\textbf{Multimodal Memory.} Extending beyond text to incorporate image, audio, and video memories with unified retrieval.

\textbf{CRDT-Based Synchronization}. Enabling seamless multi-device memory sharing without manual merging.

\subsection{Availability}

MemSt is open-source under the Apache 2.0 license:
\begin{center}
\url{https://github.com/yfyang86/memst}
\end{center}

The repository includes:
\begin{itemize}
    \item Complete Rust workspace with all crates
    \item Python bindings and FastAPI server
    \item React Web UI with KG visualization
    \item Comprehensive benchmark suite (LOCOMO)
    \item Docker deployment configurations
\end{itemize}

\vspace{1em}
\noindent\textbf{Acknowledgments}. We thank the open-source community for contributions to dependencies including PyO3, FastAPI, and HNSW. The Knowledge Graph ontology definitions were adapted from public intelligence domain taxonomies.
